% !TEX program = xelatex
% Intelligent Home Control System — Progress Report
% Compile: xelatex progress_report.tex (twice for TOC)

\documentclass[10pt,aspectratio=169]{beamer}

% ---- Theme ----
\usetheme{metropolis}
\metroset{progressbar=frametitle, numbering=fraction, sectionpage=progressbar}

% ---- Chinese (xeCJK + macOS system fonts) ----
\usepackage{fontspec}
\usepackage{xeCJK}
% 標楷體 (DFKai-SB) has no bold variant — fake bold for \textbf
\setCJKmainfont[AutoFakeBold=2.5]{標楷體}
\setCJKsansfont[AutoFakeBold=2.5]{標楷體}
\setCJKmonofont{標楷體}
\XeTeXlinebreaklocale "zh"
\XeTeXlinebreakskip = 0pt plus 1pt

% ---- Packages ----
\usepackage{booktabs}
\usepackage{tikz}
\usetikzlibrary{positioning, arrows.meta, shapes.geometric, fit, backgrounds}
\usepackage{appendixnumberbeamer}

% ---- Colors (accent for in-progress / planned) ----
\definecolor{accent}{HTML}{E07A1F}      % orange — in progress
\definecolor{planned}{HTML}{B0B0B0}     % gray — planned
\definecolor{done}{HTML}{2E8B57}        % green — done

% ---- Metadata ----
\title{智能家庭控制系統}
\subtitle{專題進度報告 — RPi Server 開發進展}
\author{趙子佾 \and 詹詠翔 \and 詹育晟}
\institute{基於 Raspberry Pi 3 與 STM32 之 BLE 通訊應用}
\date{2026 年 5 月}

% ---- Helper macros ----
\newcommand{\dn}[1]{\textcolor{done}{\textbf{#1}}}
\newcommand{\ip}[1]{\textcolor{accent}{\textbf{#1}}}
\newcommand{\pl}[1]{\textcolor{planned}{\textbf{#1}}}

\begin{document}

% ============================================================
\maketitle

% ============================================================
\begin{frame}{Outline}
  \tableofcontents
\end{frame}

% ============================================================
\section{系統架構}

\begin{frame}{硬體部署：一對多分散式節點}
  \begin{center}
  \begin{tikzpicture}[
      node distance=1.2cm and 1.6cm,
      every node/.style={font=\small},
      box/.style={draw, rounded corners=2pt, minimum width=2.2cm,
                  minimum height=0.9cm, align=center, thick},
      rpi/.style={box, fill=done!15, draw=done},
      stm/.style={box, fill=accent!15, draw=accent},
      user/.style={box, fill=blue!10, draw=blue!60},
    ]
    \node[user] (user) {使用者\\ \scriptsize (Browser)};
    \node[rpi, right=2.4cm of user] (rpi) {Raspberry Pi 3\\ \scriptsize 中央伺服器};
    \node[stm, above right=0.3cm and 2.2cm of rpi] (s1) {STM32\\ \scriptsize 客廳};
    \node[stm, right=2.2cm of rpi] (s2) {STM32\\ \scriptsize 廚房};
    \node[stm, below right=0.3cm and 2.2cm of rpi] (s3) {STM32\\ \scriptsize 房間};

    \draw[<->, thick] (user) -- node[above]{\scriptsize HTTP/WS} (rpi);
    \draw[<->, thick, accent] (rpi) -- node[above, sloped]{\scriptsize BLE} (s1);
    \draw[<->, thick, accent] (rpi) -- node[above]{\scriptsize BLE} (s2);
    \draw[<->, thick, accent] (rpi) -- node[below, sloped]{\scriptsize BLE} (s3);
  \end{tikzpicture}
  \end{center}

  \begin{itemize}
    \item RPi 同時管理多個不同空間的 STM32 節點
    \item 使用者端僅需與 RPi 互動
  \end{itemize}
\end{frame}

\begin{frame}{BLE 通訊模型}
  \begin{block}{三階段通訊流程}
    \textbf{廣播} (Advertising) $\to$
    \textbf{連線} (GATT Connect) $\to$
    \textbf{推播} (Notify)
  \end{block}

  \vspace{0.5em}
  \begin{table}
    \centering
    \small
    \begin{tabular}{@{}lll@{}}
      \toprule
      \textbf{頻道類型} & \textbf{BLE 行為} & \textbf{應用情境} \\
      \midrule
      Controller & Server $\to$ Write   & 燈光開關、溫度設定 \\
      Display    & Server $\leftarrow$ Read/Notify & 溫濕度、晃動、聲音偵測 \\
      \bottomrule
    \end{tabular}
  \end{table}
\end{frame}

% ============================================================
\section{RPi Server 設計}

\begin{frame}{技術棧}
  \begin{columns}[T]
    \column{0.5\textwidth}
    \textbf{執行環境}
    \begin{itemize}
      \item Python 3.11+（dev 用 3.12）
      \item uv \, 套件管理
      \item Taskfile \, 統一指令入口
    \end{itemize}

    \textbf{Web 層}
    \begin{itemize}
      \item Flask + Flask-Login
      \item Flask-WTF（CSRF）
      \item Jinja2 模板
    \end{itemize}

    \column{0.5\textwidth}
    \textbf{資料層}
    \begin{itemize}
      \item SQLite + WAL
      \item 自製 repository 層
    \end{itemize}

    \textbf{BLE 層}
    \begin{itemize}
      \item bluepy（Linux/RPi）
      \item Mock 實作（Windows dev）
    \end{itemize}

    \textbf{品質工具}
    \begin{itemize}
      \item ruff、mypy strict
      \item pytest（103 tests）
    \end{itemize}
  \end{columns}
\end{frame}

\begin{frame}{模組分層}
  \begin{center}
  \begin{tikzpicture}[
      every node/.style={font=\small},
      layer/.style={draw, rounded corners=2pt, minimum width=8cm,
                    minimum height=0.85cm, align=center, thick},
    ]
    \node[layer, fill=blue!10]   (web) {\textbf{web/} \, Flask blueprints、auth、templates};
    \node[layer, fill=done!15, below=0.25cm of web]   (svc) {\textbf{services/} \, user / device / channel \, 業務邏輯};
    \node[layer, fill=orange!15, below=0.25cm of svc] (db)  {\textbf{db/} \, repository（users / devices / channels / readings）};
    \node[layer, fill=accent!15, below=0.25cm of db]  (ble) {\textbf{ble/} \, interface / parser / rate\_limiter / managers};
    \node[layer, fill=gray!15, below=0.25cm of ble]   (core){\textbf{core/} \, config、logging};

    \draw[-{Latex[length=2mm]}, thick] (web) -- (svc);
    \draw[-{Latex[length=2mm]}, thick] (svc) -- (db);
    \draw[-{Latex[length=2mm]}, thick] (svc) -- ++(3.5,-0.55) |- (ble);
  \end{tikzpicture}
  \end{center}

  \begin{itemize}
    \item 上層只依賴下層；service 層為純業務邏輯
    \item BLE 介面用 \texttt{Protocol} 抽象，方便 mock 測試
  \end{itemize}
\end{frame}

% ============================================================
\section{已完成進度}

\begin{frame}{\dn{已完成} Phase 1–2：骨架與 BLE 通訊鏈結}
  \textbf{Phase 1：骨架}
  \begin{itemize}
    \item uv 專案、目錄結構、\texttt{config.py} 讀環境變數
    \item \texttt{db/schema.sql}：users / devices / channels / readings + WAL
    \item \texttt{\_\_main\_\_.py}：可啟動 Flask app，含 \texttt{/health}
  \end{itemize}

  \vspace{0.5em}
  \textbf{Phase 2：BLE 通訊}
  \begin{itemize}
    \item \texttt{ble/interface.py}：\texttt{BLEManager} Protocol 抽象
    \item \texttt{ble/parser.py}：10 種格式雙向轉換（int/uint 8/16/32、float32，LE/BE）
    \item \texttt{ble/rate\_limiter.py}：per-key 限頻，可注入 fake clock
    \item \texttt{ble/mock\_manager.py} / \texttt{bluepy\_manager.py}：Windows dev 與 RPi 雙實作
    \item \texttt{bluepy\_manager}：每個 peripheral 一個 worker thread + command queue，避開 bluepy 跨執行緒問題
  \end{itemize}
\end{frame}

\begin{frame}{\dn{已完成} Phase 3a：DB Repository 層}
  四個 repository 模組，每個自帶領域錯誤型別：
  \begin{itemize}
    \item \texttt{db/users.py}：\texttt{create / get\_by\_id / get\_by\_username / update\_password}
    \item \texttt{db/devices.py}：\texttt{create / get\_by\_*  / list\_* / delete}
    \item \texttt{db/channels.py}：含 \texttt{Literal["controller", "display"]} 型別
    \item \texttt{db/readings.py}：\texttt{insert / list\_by\_channel / delete\_older\_than}
  \end{itemize}

  \vspace{0.5em}
  \begin{block}{設計重點}
    \begin{itemize}
      \item 時間戳統一 UTC \texttt{"YYYY-MM-DD HH:MM:SS"} 字串
      \item 領域錯誤：\texttt{DuplicateUsernameError}、\texttt{DeviceNotFoundError} …
      \item 測試 fixture：in-memory SQLite + schema 已套用
    \end{itemize}
  \end{block}
\end{frame}

\begin{frame}{\dn{已完成} Phase 3b：Service 層}
  \textbf{UserService}
  \begin{itemize}
    \item bcrypt 雜湊（cost 可注入，測試用低 cost）
    \item 密碼長度 8–72 bytes（避免 bcrypt 截斷）
  \end{itemize}

  \textbf{DeviceService}（建構注入 \texttt{BLEManager}）
  \begin{itemize}
    \item \texttt{scan / add\_device / remove\_device / list\_devices}
    \item \texttt{add\_device}：驗證 MAC → 寫 DB → best-effort 連線
  \end{itemize}

  \textbf{ChannelService}（注入 BLE、RateLimiter、\texttt{on\_reading} callback）
  \begin{itemize}
    \item \texttt{write\_command}：僅 controller 型，編碼後 BLE write
    \item \texttt{handle\_notify}：解析 → 即時推播（不限頻）→ DB 寫入（限頻）
    \item \texttt{get\_history / list\_by\_device}
  \end{itemize}
\end{frame}

\begin{frame}{\dn{已完成} Phase 3c：認證 Blueprint}
  \begin{itemize}
    \item \texttt{web/\_\_init\_\_.py}：application factory \texttt{create\_app}
    \item 整合 \textbf{Flask-Login} + \textbf{Flask-WTF CSRF}
    \item \texttt{web/db.py}：per-request SQLite 連線（\texttt{flask.g} + teardown）
    \item \texttt{web/auth.py}：\texttt{/auth/register}、\texttt{/auth/login}、\texttt{/auth/logout}
    \item \texttt{next} 參數的 \textbf{open-redirect 防護}
    \item Jinja2 陽春模板（base / index / auth）含 CSRF token
  \end{itemize}
\end{frame}

% ============================================================
\section{進行中與規劃}

\begin{frame}{\ip{進行中} Phase 3d：Device / Channel CRUD}
  尚待補齊的 Blueprint：
  \begin{itemize}
    \item \texttt{/devices}：列表、掃描、新增、刪除
    \item \texttt{/devices/<id>/channels}：新增、刪除
    \item \texttt{POST /channels/<id>/write}：控制型頻道下指令
  \end{itemize}

  \vspace{0.6em}
  \textbf{設計重點}
  \begin{itemize}
    \item 重用既有 \texttt{DeviceService} / \texttt{ChannelService}
    \item Server-rendered 表單 + CSRF；後續再以 HTMX 強化體驗
  \end{itemize}
\end{frame}

\begin{frame}{\pl{規劃中} Phase 3e：SocketIO 即時推播與前端}
  \textbf{後端}
  \begin{itemize}
    \item Flask-SocketIO（threading 模式）
    \item 與 bluepy worker thread 串接
    \item 每個監控型頻道一個 SocketIO room
  \end{itemize}

  \textbf{前端}
  \begin{itemize}
    \item Jinja2 + \textbf{HTMX} 漸進增強
    \item \textbf{Chart.js} 繪製歷史趨勢圖
    \item 訂閱 Notify 事件後即時更新圖表
  \end{itemize}
\end{frame}

\begin{frame}{\pl{規劃中} STM32 韌體開發}
  \textbf{硬體：B-L475E-IOT01A}
  \begin{itemize}
    \item HTS221 \, 溫濕度感測器
    \item LSM6DSL \, 加速度 / 陀螺儀
    \item MP34DT01 \, 數位麥克風
    \item NFC、使用者 LED
  \end{itemize}

  \textbf{韌體任務}
  \begin{itemize}
    \item 開發 BLE GATT Service \, /\, Characteristic 設計
    \item 感測器資料採集與封包格式（對齊 RPi 端 parser）
    \item 接收控制指令（LED、設定值）
    \item 工具鏈：\textbf{STM32CubeIDE} + ST HAL
  \end{itemize}
\end{frame}

\begin{frame}{\pl{規劃中} 端到端整合}
  \begin{itemize}
    \item 多 STM32 節點\textbf{分散式佈署測試}（客廳 / 廚房 / 房間）
    \item RPi $\leftrightarrow$ STM32 \textbf{真機聯調}
    \item BLE \textbf{Notify 推播延遲}量測
    \item 異常情境：斷線重連、信號弱、節點重啟
  \end{itemize}
\end{frame}

\begin{frame}{\pl{規劃中} Phase 4：部署與穩定性}
  \begin{itemize}
    \item \texttt{systemd} service 安裝腳本（boot 即啟動）
    \item RPi 上長時間穩定性測試（記憶體 / 連線 / DB 大小）
    \item 部署文件與 troubleshooting guide
    \item 成果展示準備
  \end{itemize}
\end{frame}

\begin{frame}[standout]
  謝謝聆聽！\\
\end{frame}

\end{document}
